\chapter{Propositional Logic}

\emph{Propositional logic}, also called Boolean logic, is one of the foundational branches of mathematics and, thus, the 
first chapter of this book. In propositional logic, we use symbolic variables to represent the logic, and we can use any symbol to represent a proposition, such as 
\(p, q, r, s, t,\) etc. Propositions can be either \emph{true} (\(1\), T) or \emph{false} (\(0\), F), but they cannot be both.

\textbf{Remark}: In programming languages like C, any number greater than 0 is treated as T.

\section{Predicates}

A \emph{predicate} is a statement which contains variables linked by some property stated in the language of formal logic, and whose value can be true or false.

\section{Predicate Substitution}

Given a predicate \(P\) on the variable \(v\), we can substitute every free instance of \(v\) with the term \(v\) to extend 
our predicate to another variable space.

\[
	[t/v] P
\]


Types of substitution:

\begin{itemize}
	
	\item \emph{Full replacement}: \([y/x]\forall z((z = x) \land (y \in x))\), converted to \(\forall z((z = y) \land (y \in y))\).

	\item \emph{Partial replacement}: \([y/ x]x = y \land \forall x (x \in y)\), the \(x\) of \(\forall x(x \in y)\) is not a free variable, hence it can not be
	replaced completely, just \(y = y \land \forall x (x \in y)\).

\end{itemize}

\textbf{Example:}

Prove \(\forall a (a = a)\):

Let \(P:= x = a\) be a predicate in \(x\) and \(a\). We apply:

\[
	[a /x](x = a) \text{ to } a = a \text{, then by extension of our predicate } \forall a (a = a).
\]

Prove \(\forall a \forall b (a = b \implies b = a )\):

Let \(P:= x = a\) be a predicate in \(a\) and \(x\). We apply:

\[
	[b / x](x = a) \text{ to } b = a \text{, then by extension of our predicate } \forall a \forall ( a = b \implies b = a )).
\]

\QED 

Prove \(\forall a (a = b \land b = c \implies c = a)\):

We know \(a = b \land b = a\), let this be our predicate with \(x\) and apply the substitution \([c/b]\), which gives 
us \(a = x\), \(a = b\) and finally \(a = c\).

\QED

\section{Logic Operators and Truth Tables}

\smallskip
\begin{multicols}{2}

	\subsection*{NOT (\(\neg\), \(\sim\))}

	\begin{tabular}{|c|c|}
		\toprule
		\(p\) & \(\neg q\) \\
		\midrule
		T   & F        \\
		F   & T        \\
		\bottomrule
	\end{tabular}

	\vspace{1em}

	\subsection*{AND (\(\land\))}

	\begin{tabular}{|c|c|c|}
		\toprule
		\(p\) & \(q\) & \(p \land q\) \\
		\midrule
		T   & T   & T           \\
		T   & F   & F           \\
		F   & T   & F           \\
		F   & F   & F           \\
		\bottomrule
	\end{tabular}

	\vspace{1em}

	\subsection*{OR (\(\lor\))}

	\begin{tabular}{|c|c|c|}
		\toprule
		\(p\) & \(q\) & \(p \lor q\) \\
		\midrule
		T   & T   & T          \\
		T   & F   & T          \\
		F   & T   & T          \\
		F   & F   & F          \\
		\bottomrule
	\end{tabular}

	\vspace{1em}

	\subsection*{IMPLIES (\(\implies\))}

	\begin{tabular}{|c|c|c|}
		\toprule
		\(p\) & \(q\) & \(p \implies q\) \\
		\midrule
		T   & T   & T        \\
		T   & F   & F        \\
		F   & T   & T        \\
		F   & F   & T        \\
		\bottomrule
	\end{tabular}

	\columnbreak

	\subsection*{IFF (\(\iff\))}

	\begin{tabular}{|c|c|c|}
		\toprule
		\(p\) & \(q\) & \(p \iff q\) \\
		\midrule
		T   & T   & T         \\
		T   & F   & F         \\
		F   & T   & F         \\
		F   & F   & T         \\
		\bottomrule
	\end{tabular}

	\vspace{1em}

	\subsection*{XOR (\(\oplus\))}

	\begin{tabular}{|c|c|c|}
		\toprule
		\(p\) & \(q\) & \(p \oplus q\) \\
		\midrule
		T   & T   & F            \\
		T   & F   & T            \\
		F   & T   & T            \\
		F   & F   & F            \\
		\bottomrule
	\end{tabular}

	\vspace{1em}

	\subsection*{NOR (\(\downarrow\))}

	\begin{tabular}{|c|c|c|}
		\toprule
		\(p\) & \(q\) & \(p \downarrow q\) \\
		\midrule
		T   & T   & F                \\
		T   & F   & F                \\
		F   & T   & F                \\
		F   & F   & T                \\
		\bottomrule
	\end{tabular}

	\vspace{1em}

	\subsection*{NAND (\(\uparrow\))}

	\begin{tabular}{|c|c|c|}
		\toprule
		\(p\) & \(q\) & \(p \uparrow q\) \\
		\midrule
		T   & T   & F              \\
		T   & F   & T              \\
		F   & T   & T              \\
		F   & F   & T              \\
		\bottomrule
	\end{tabular}

\end{multicols}
\medskip

\section{Implications}

As a side note, for the \emph{implication} 

\[
	p \implies q
\]

we can think about it as \emph{if, then}. For the cases  false, true; and false, false; we call them \emph{vacuously} true. 
From a false premise, anything can follow.

\textbf{Example:}

Let \(p\) be ``it rained'' and \(q\) ``the street is wet'', with theirs respective negations ``it did not rain'', ``the street is not wet''. 

Now, let us study each of the cases:

\begin{center}
		\begin{tabular}{|c|c|c|}
		\toprule
		\(p\) & \(q\) & \(p \implies q\) \\
		\midrule
		T   & T   & T        \\
		T   & F   & F        \\
		F   & T   & T        \\
		F   & F   & T        \\
		\bottomrule
	\end{tabular}
\end{center}

\begin{itemize}

	\item \(p\) and \(q\) are both true: ``It rained and the street is wet.'' This makes the implication \(p \implies q\) true.

	\item \(p\) is true and \(q\) is false: ``It rained but the street is not wet.'' This makes the implication \(p \implies q\) false, since the stree should be wet.

	\item \(p\) is false and \(q\) is true: ``It did not rain but the street is wet.'' This makes the implication \(p \implies q\) true (vacuously true). 
	An external factor could have been responsible for the street being wet.

	\item \(p\) and \(q\) are both false: ``It did not rain and the street is not wet.'' This makes the implication \(p \implies q\) true (vacuously true).

\end{itemize}

We simplify the possibilites inside this mini universe and the implication becomes very intuitive.

\section{If And Only If}

\emph{If and only if} \(\iff\) is equivalent to  

\[
	(A \implies B) \land (B \implies A).
\]

Which is justs its definition.

\section{Tautology and Contradiction}

\emph{Tautology}: A logical formula that is always true.

Common tautologies 

\begin{itemize}

	\item \(p \implies q \lor p\)

	\item \(q \implies p \lor q\)
	
	\item \(p \land q \implies p\)
	
	\item \(q \land q \implies q \)

	\item \([(p \implies q) \land (q \implies r)] \implies  (p \implies r)\)

\end{itemize}

\emph{Contradiction}: A formula that is always false.

\section{Logical Equivalences}

Two logical formulas \(P, Q\) are said to be \emph{logically equivalent} if they have the same truth value in every possible scenario.
We can write this as \(P \equiv Q\) or \(P \iff Q\).

\subsection{Laws of Logical Equivalences}

\emph{Commutative Laws}

\[
	p \land q \iff q \land p \qquad p \lor q \iff q \lor p
\]

\emph{Associative Laws}

\[
	(p \land q) \land r \iff p \land (q \land r) \qquad (p \lor q) \lor r \iff p \lor (q \lor r)
\]

\emph{Distributive Laws}

\[
	p \land (q \lor r) \iff (p \land q) \lor (p \land r) \qquad
	p \lor (q \land r) \iff (p \lor q) \land (p \lor r)
\]

\emph{Identity Laws}

\[
	p \land \text{T} \iff p \qquad p \lor \text{F} \iff p
\]

\emph{Negation Laws}

\[
	p \lor \sim p \iff \text{T} \qquad p \land \sim p \iff \text{F}
\]

\emph{Double Negation Law}

\[
	\sim(\sim p) \iff p
\]

\emph{Idempotent Laws}

\[
	p \land p \iff p \qquad p \lor p \iff p
\]

\emph{Universal Bound Laws}

\[
	p \lor \text{T} \iff \text{T} \qquad p \land \text{F} \iff \text{F}
\]

\emph{De Morgan’s Laws}

\[
	\sim (p \land q) \iff (\sim p) \lor (\sim q) \qquad
	\sim (p \lor q) \iff (\sim p) \land (\sim q)
\]

\emph{Absorption Laws}

\[
	p \lor (p \land q) \iff p \qquad p \land (p \lor q) \iff p
\]

\emph{Conditional Laws}

\[
	(p \implies q) \iff (\sim p \lor q) \qquad \sim(p \implies q) \iff (p \land \sim q)
\]

\emph{Complement Law}

\[
	p \lor \neg p \iff \text{T} \qquad p \land \neg p \iff \text{F}
\]

\emph{Biconditional Law}

\[
	p \iff q \iff (p \implies q) \land (q \implies p)
\]

\emph{Transitivity Law}

\[
	(p \implies q) \land (q \implies r) \implies (p \Rightarrow r)
\]

\emph{Indirect Proof (Contrapositive)}

\[
	(p \implies q) \iff (\neg q \implies \neg p)
\]

\emph{Disjunctive Syllogism (Disjunctive Exclusion)}

\[
	p \nabla q \equiv (p \lor q) \land \neg p \implies q
\]

\[
	p \nabla q \equiv(p \land q) \lor \neg  (p \land q)
\]

\section{Truth Tables}

\emph{Truth tables} are a fundamental tool in logic that systematically show the truth value
(true or false) of a compound statement for every possible combination of the truth values of 
its individual component statements.

Essentially, they lay out all the scenarios and the resulting truth of the overall logical expression.
This helps determine if an argument is valid, if statements are logically equivalent, or the circumstances
under which a complex statement is true or false.

\textbf{Example:}

\begin{center}
	\begin{tabular}{|c|c|c|c|c|c|c|c|}
		\hline
		\(p\) & \(q\) & \(\neg p\) & \(\neg q\) & \(\neg p \implies q\) & \((\neg p \implies q) \land \neg p\) & \(\left[(\neg p \implies q) \land \neg p\right] \Rightarrow q\) \\
		\hline
		T   & T   & F        & F        & T                      & F                                     & T                                                                \\
		T   & F   & F        & T        & T                      & F                                     & T                                                                \\
		F   & T   & T        & F        & T                      & T                                     & T                                                                \\
		F   & F   & T        & T        & F                      & F                                     & T                                                                \\
		\hline
	\end{tabular}
\end{center}

\subsection{Filling a Truth Table}
To fill a truth table for a logical expression with truth values (True or False), you follow a 
specific order for the input variables. This order ensures that all possible combinations of truth values 
for the variables are covered.

\begin{enumerate}

	\item \textbf{List all possible combinations of truth values for the input variables}: 
			If you have \(n\) variables, the number of rows in the truth table will be \(2^n\). Each 
			variable can be either True (T) or False (F).

	\item \textbf{Order of the input variables}:

			\begin{itemize}
					
				\item Start by filling in the truth values for the first variable. It alternates between 
					  True and False every \(2^{n-1}\) rows.
					
				\item Then for the second variable, it alternates every \(2^{n-2}\) rows, and so on.
					
				\item In short: the first variable alternates every other row, the second variable every two 
					rows, the third every four rows, and so on.
			
			\end{itemize}

\end{enumerate}

\textbf{Example:} 

For 3 variables, there are \(2^3 = 8\) possible combinations of truth values. The truth values are filled 
in the following order:

\[
	\begin{array}{|c|c|c|c|}
		\hline
		p & q & r & \text{Expression Result} \\
		\hline
		T & T & T &       -                  \\
		T & T & F &       -                  \\
		T & F & T &       -                  \\
		T & F & F &       -                  \\
		F & T & T &       -                  \\
		F & T & F &       -                  \\
		F & F & T &       -                  \\
		F & F & F &       -                  \\
		\hline
	\end{array}
\]

The pattern for filling the input truth values

\begin{itemize}

	\item The first column \(p\) alternates every 4 rows: `T, T, F, F, T, T, F, F`.

	\item The second column \(q\) alternates every 2 rows: `T, T, F, F, T, T, F, F`.

	\item The third column \(r\) alternates every row: `T, F, T, F, T, F, T, F`.

\end{itemize}

This ensures that all combinations of \(p\), \(q\), and \(r\) are covered, and you 
can then evaluate the logical expression for each combination.

\subsubsection{Truth Table for the Expression \texorpdfstring{\( (p \land q) \lor q \)}{}}

\[
	\begin{array}{|c|c|c|c|}
		\hline
		p & q & r & (p \land q) \lor r \\
		\hline
		T & T & T & T                  \\
		T & T & F & T                  \\
		T & F & T & T                  \\
		T & F & F & F                  \\
		F & T & T & T                  \\
		F & T & F & F                  \\
		F & F & T & T                  \\
		F & F & F & F                  \\
		\hline
	\end{array}
\]

\subsection{Disjunctive Normal Form (DNF)}

\emph{Disjunctive Normal Form (DNF)} is a standard way of writing a logical expression as a 
\emph{disjunction(OR)} of \emph{conjunctions (ANDs)}. A DNF expression consists of a series of
conjunctions of literals, where each conjunction is connected by disjunctions.

To find the DNF in a truth table take the rows of the final result where there are
true statements and bind the propositions that generated it with an AND inside parenthesis.
Repeat it with each of true rows and connect all parenthesis with OR's.

\textbf{Example:}

Consider the logical expression:

\[
	(p \land q) \lor (\neg p \land r) \lor (q \land \neg r).
\]

This is in DNF because it is a disjunction (OR) of conjunctions (ANDs) of literals.

\subsection{Conjunctive Normal Form (CNF)}

\emph{Conjunctive Normal Form (CNF)} is a standard way of writing a logical expression as a 
\emph{conjunction (AND)} of \emph{disjunctions (ORs)}. A CNF expression consists of a series of 
disjunctions of literals, where each disjunction is connected by conjunctions.

To find the CNF proceed just as the DNF but with the false rows and instead
of ANDs inside the parenthesis use OR and connect the terms with OR\@. 
Also add a negation before each parenthesis.

\textbf{Example:}

Consider the logical expression:

\[
	\neg (p \lor q) \land \neg (\neg p \lor r) \land \neg (q \lor \neg r).
\]

This is in CNF because it is a conjunction (AND) of disjunctions (ORs) of literals.

\section{Basis}

A \emph{Basis} in the field of logic is a set of operations with which you can 
represent all possible logic operations.

\textbf{Example:} The set \(\{ \lnot, \land\}\)

\section{Karnaugh Maps}

\emph{Karnaugh Maps (K-Maps)} are a graphical method used to simplify Boolean expressions. 
The main goal of a K-map is to group adjacent cells that contain 1's in order to simplify the expression. 
A K-map helps identify common terms, allowing the Boolean expression to be reduced to its simplest form.

\subsection{Karnaugh Map for Two Variables}

Consider the Boolean expression \( (A \lor (B \land \neg A \land \neg B)) \).

We first construct a K-map for two variables, \(A\) and \(B\). The truth table for this expression 
gives the following values:

\[
	\begin{array}{|c|c|c|}
		\hline
		A & B & (A \lor (B \land \neg A \land \neg B)) \\
		\hline
		0 & 0 & 0                                      \\
		0 & 1 & 1                                      \\
		1 & 0 & 1                                      \\
		1 & 1 & 1                                      \\
		\hline
	\end{array}
\]

The corresponding K-map is:

\[
	\begin{array}{|c|c|c|c|c|}
		\hline
		AB           & 00 & 01 & 11 & 10 \\
		\hline
		\text{Value} & 0  & 1  & 1  & 1  \\
		\hline
	\end{array}
\]

Here, we group the ones together to simplify the Boolean expression. The simplified expression is:

\[
	A \lor B.
\]

\subsection{Karnaugh Map for Three Variables}

Now, let's consider the expression \( \neg C \). This expression only 
depends on one variable, but for illustration, we will use a 
3-variable K-map with variables \(A\), \(B\), and \(C\).

The truth table for \( \neg C \) is as follows:

\[
	\begin{array}{|c|c|c|c|}
		\hline
		A & B & C & \neg C \\
		\hline
		0 & 0 & 0 & 1      \\
		0 & 0 & 1 & 0      \\
		0 & 1 & 0 & 1      \\
		0 & 1 & 1 & 0      \\
		1 & 0 & 0 & 1      \\
		1 & 0 & 1 & 0      \\
		1 & 1 & 0 & 1      \\
		1 & 1 & 1 & 0      \\
		\hline
	\end{array}
\]

The corresponding K-map for three variables \(A\), \(B\), and \(C\) is:

\[
	\begin{array}{|c|c|c|c|}
		\hline
		AB \backslash C & 0 & 1 \\
		\hline
		00              & 1 & 0 \\
		01              & 1 & 0 \\
		11              & 1 & 0 \\
		10              & 1 & 0 \\
		\hline
	\end{array}
\]

We see that the ones are grouped in a column, leading to the simplified Boolean expression:

\[
	\neg C.
\]

\subsection{Solving a Karnaugh Map}

To solve a Karnaugh Map and simplify a Boolean expression, follow these steps:

\begin{enumerate}
	
	\item \textbf{Determine the Number of Variables:} Decide how many variables are in the 
		   Boolean function. This determines the size of the K-Map:
	      
		\begin{itemize}
		 
			\item 2 variables: \(2 \times 2\).
		 	
			\item 3 variables: \(2 \times 4\).
		 
			\item 4 variables: \(4 \times 4\).
		 
			\item etc.
	  
		\end{itemize}

	\item \textbf{Fill in the K-Map:} Place 1's in the cells that correspond to the min-terms 
		  (where the function outputs 1). You may also include do not care conditions 
		  (usually denoted as \(X\)).

	\item \textbf{Group the 1's:} Form groups (called \emph{implicants}) of 1's. 
		   The groups must follow these rules:
	      
		\begin{itemize}
		 
			\item Each group must contain \(1, 2, 4, 8, \ldots\) (powers of 2) 1's.
		 
			\item Groups must be rectangular (e.g., \(1 \times 2\), \(2 \times 2\)).
		 
			\item Groups can wrap around the edges of the K-Map.
		 
			\item Try to form the largest groups possible to simplify the expression.
		 
			\item Each 1 should be included in at least one group.
	 
		\end{itemize}

	\item \textbf{Write the Simplified Expression:} For each group:

		  \begin{itemize}

			\item Identify the variables that are constant (either always 0 or always 1) across the group.

			\item Write a product term (AND) using only the constant variables.

			\item Combine all product terms with OR operations to get the final simplified SOP (Sum of 
				  Products) expression.

		\end{itemize}

	\end{enumerate}

\section{Open Sentences}

When writting statements we normally do not write the statement with concrete values, but instead define it in terms of variables like 
\(x\) which is used to denote the members of a collection for which our statement holds or not holds. We denote this as 
an \emph{open sentence}.

\textbf{Example:}

\[
	P(x): \text{ The sum of the angles of } x \text{ is } 180^{\circ}, \, \forall x \in S
\]

where \(S\) is the set of the triangles for which \(P\) holds.

\section{Mathematical Quantifiers}

\emph{Mathematical quantifiers} are symbols used as shorthand for common quantifying of elements when writting mathmatical statements.

\subsection{Universal Quantifier \texorpdfstring{\(\forall\)}{}}  

Means \textbf{``for all''} or \textbf{``for every''}.

\emph{Negation:} \(\neg (\forall x)\,P(x) \equiv (\exists x)\, \neg P(x)\) ``Not all'' is the same as 
``There exists one that does not''.

\[
	\neg (\forall x, P(x)) \equiv \exists x, \neg P(x)
\]

\textbf{Proof:}

We want to prove that \(\neg (\forall i \in I : E(i)) \iff (\exists i \in I: \neg E(i))\), given \(I \neq \emptyset\)

\begin{align*}
	\neg (\forall i \in I : E(i)) &\iff \neg (I = \{i \in I \mid E(i)\})\\
								  &\iff I \neq \{i \in I \mid E(i)\} \\
								  &\iff \{i \in I \mid \neg E(i)\} \neq \emptyset \\
								  &\iff (\exists i \in I : \neg E(i)
\end{align*}

\QED

\textbf{Example:} 

\(\forall x \in \R,\ x^2 \geq 0\) For all real numbers, the square is greater than or equal to zero.

\textbf{Negated Example:} 

\(\exists x \in \R,\ x^2 < 0\) There exists a real number whose square is less than zero — this is 
false.

\subsection{Existential Quantifier \texorpdfstring{\(\exists\)}{}}  
	
Means \textbf{``there exists at least one''}.
 		
\emph{Negation:}  \(\neg (\exists x)\,P(x) \equiv (\forall x)\, \neg P(x)\) ``There does not exist'' is 
the same as ``For all, not''.

\textbf{Proof:}

We want to prove that \(\neg (\exists i \in I : E(i)) \iff (\forall i \in I: \neg E(i))\), given \(I \neq \emptyset\)

\begin{align*}
	\neg (\exists i \in I : E(i)) &\iff \neg (I = \{i \in I \mid E(i)\} \neq \emptyset)\\
								  &\iff  I \neq \{i \in I \mid E(i)\} = \emptyset \\
								  & \iff \forall i \in I: \neg E(i)
\end{align*}

\QED

\textbf{Example:} 

\(\exists x \in \N,\ x > 10\) There exists a natural number greater than 10.
			
\textbf{Negated Example:}  

\(\forall x \in \N,\ x \leq 10\) All natural numbers are less than or equal to 10 — this is false.

\subsection{Unique Existential Quantifier \texorpdfstring{\(\exists ! \)}{}} 
	
Means \textbf{``there exists exactly one''}.
		
\emph{Negation:} ``Not exactly one'' means 

\[ 
	\neg (\exists! x)\, P(x) \equiv (\forall x)\, \neg P(x)\ \lor\ (\exists x_1 \neq x_2)\, P(x_1) \land  P(x_2)
\] 

Either no such \(x\) exists, or more than one does.

\textbf{Example:} 

\(\exists! x \in \R,\ x + 5 = 0\) There exists exactly one real number such that \( x + 5 = 0 \).

\textbf{Negated Example:}  

\(\exists x_1 \neq x_2 \in \R,\ x_1^2 = 4 \land x_2^2 = 4\) There are multiple solutions to 
\( x^2 = 4 \).

\subsection{Negations Remark}

When dealing with proofs, it is not uncommon that we will need to negate an statement for the purpose of prooving 
a simpler equivalent statement. For example: 

Note that \(p \implies q\) is logically equivalent to \(\neg p \lor q\). Hence, if would to prove its negation we can use the following equivalences:

\[
	\neg (p \implies q) \equiv \neg (\neg p \lor q) \equiv p \land \neg q.
\]


\subsection{Order Remark}

When using quantifiers, \emph{the order matters}, it is not same to write 

\[
	\forall x \in M, \exists y \in N: \, E(x, y)
\]

and 

\[
	\exists y \in N, \forall x \in M: \, E(x, y)
\]

Note that in the first example \(y\) can depend on \(x\) which means that for each \(x\) the may be different \(y\); while in the second example 
we declare one or more \(y\)'s which satisfy the condition for all \(x\).


\section{Equivalence Reformulations}

An \emph{equivalence reformulation} is when we trasnform a statement from one form to another which is \emph{logically equivalent}, it has the values in a truth table. 

\[
	A \iff B
\]

We use them all the time during our proofs as the connectors when we want to rewrite a statement is form which more easier to work with. 

\textbf{Example:}

\[
	|x| \le y \iff (x \le y \land -x \le y) \iff -y \le x \le y
\]



